\subsection{Planteamiento del Problema}

Las organizaciones que atienden consultas por canales de mensajería, como universidades, empresas, tiendas, centros de servicios y entidades públicas, enfrentan dificultades recurrentes en su comunicación institucional digital con estudiantes, clientes, ciudadanos o beneficiarios.

En este entorno, estas instituciones y negocios reciben un alto volumen de consultas repetitivas sobre horarios, requisitos, estados de trámite, precios, disponibilidad, costos, fechas clave y soporte operativo. La gestión manual de estas interacciones suele generar demoras en la respuesta, saturación del personal, inconsistencias en la información y baja trazabilidad del historial conversacional.

El problema se intensifica cuando existen múltiples canales no integrados, porque se duplican esfuerzos y se dificulta mantener criterios unificados de atención. Además, la falta de automatización limita la capacidad de ofrecer respuestas oportunas fuera del horario laboral, afectando la experiencia del usuario y la percepción de calidad del servicio.

En consecuencia, el problema central consiste en la ausencia de una solución \textit{Software como Servicio (SaaS)} que centralice conversaciones, automatice respuestas frecuentes y optimice los flujos de comunicación con control, escalabilidad y consistencia informativa.
\subsubsection{Formulación del Problema}
¿Cómo mejorar la atención y gestión de consultas de los usuarios mediante la implementación de un sistema automatizado de respuestas integrado en plataformas de mensajería como WhatsApp y Telegram?